\subsection{Acercamiento a la solución}

\textbf{Contexto del desarrollo y justificación.} Carer-Record complementa la FCE, que continúa siendo el registro oficial. La solución reúne la consulta de información autorizada de la FCE con notas de contexto propias, para apoyar la continuidad informativa durante la hospitalización. No reemplaza la FCE ni escribe en ella.

\textbf{Solución técnica.} La arquitectura del producto final se sustenta en las siguientes decisiones:

\begin{itemize}
    \item \textbf{Aplicación móvil:} Flutter permite construir una interfaz orientada a teléfonos del personal asistencial, con acceso rápido y baja fricción cognitiva.
    
    \item \textbf{Servicios de aplicación:} Supabase Auth, PostgreSQL con seguridad a nivel de fila (Row-Level Security, RLS), PostgREST, llamadas a procedimientos remotos (Remote Procedure Call, RPC) y Edge Functions proveen identidad, persistencia y operaciones autorizadas.
    
    \item \textbf{Integración hospitalaria:} la consulta a la FCE se realiza mediante una pasarela autorizada; la aplicación móvil no consulta directamente la fuente hospitalaria.

    \item \textbf{Procesamiento local:} el ASR y la asistencia de redacción y resumen se ejecutan en el dispositivo. El audio y la transcripción sin confirmar permanecen allí; solo la nota confirmada se persiste en el backend.
\end{itemize}

\textbf{Delimitación académica.} El producto consulta fuentes autorizadas; la demostración utiliza exclusivamente datos sintéticos y no equivale a conexión o validación productiva hospitalaria. El contrato de cada integración se define con su fuente; la fuente de demostración descrita utiliza Fast Healthcare Interoperability Resources (FHIR), versión 4 (R4).

\begin{itemize}
    \item \textbf{Entorno de demostración:} se utiliza una FCE simulada con FHIR R4 en Docker y datos sintéticos.
    \item \textbf{Privacidad y datos:} la demostración académica emplea datos sintéticos y no información clínica real.
    \item \textbf{Alcance de usuarios:} el sistema se orienta exclusivamente al personal asistencial; la comunicación directa con familiares queda fuera del prototipo.
\end{itemize}
